\scanpage{100}

\mspage{VII - 9 -}

n\textsuperscript{o}339

tifs des éléments de $P(A)$.

Ainsi, pour $x\handwritten{\in}K^{\handwritten{*}}$, on définit des entiers rationnels $v_P(x)$ en posant

\begingroup\csname tagsleft@true\endcsname
\[ \operatorname{div}(x)=\handwritten{\sum}_{P\handwritten{\in}P(A)}v_P(x)\mathbin{.}P\ .\tag*{(3)} \]
\endgroup

Des relations $\operatorname{div}(xy)=\operatorname{div}(x)+\operatorname{div}(y)$ et $\operatorname{div}(x+y)\handwritten{\geq}\inf(\operatorname{div}(x),\operatorname{div}(y))$, on déduit que les $v_P$ sont des valuations discrètes. Pour que $x\handwritten{\in}A$, il faut et il suffit que $\operatorname{div}(x)\handwritten{\geq}0$, c'est à dire que $v_P(x)\handwritten{\geq}0$ pour tout $P\handwritten{\in}P(A)$. Ainsi les $v_P$ vérifient les conditions (AK.1) et (AK.2), et évidemment aussi (AK.3).

C.Q.F.D.

COROLLAIRE.- \uline{Pour qu'un anneau noethérien soit un anneau de Krull}, \uline{il faut et il suffit qu'il soit intégralement clos}.

En effet un anneau noethérien intégralement clos est complètement intégralement clos (Chap.V,\S1,n\textsuperscript{o}3).

4.- \uline{Valuations essentielles d'un anneau de Krull}.

Soit $A$ un anneau de Krull. Les éléments de $D(A)$ sont alors appelés les \uline{diviseurs} de $A$, et les classes de $D(A)$ suivant le sous-groupe des diviseurs principaux sont appelées les \uline{classes de diviseurs} de $A$. On appelle \uline{valuations essentielles} de $A$ les valuations $v_P$ définies, comme au n\textsuperscript{o} précédent, par la formule

\begingroup\csname tagsleft@true\endcsname
\[ \operatorname{div}(x)=\handwritten{\sum}_{P\handwritten{\in}P(A)}v_P(x)\mathbin{.}P\tag*{(4)} \]
\endgroup

pour $x\handwritten{\in} K^{\handwritten{*}}$. On a remarqué, au cours de la démonstration du th.1, que les valuations $v_P$ vérifient les propriétés (AK.1), (AK.2) et (AK.3) de la déf.2.
En outre, ces valuations discrètes $v_P$ sont \uline{normées} : en effet, pour tout diviseur extrémal $P\handwritten{\in}P(A)$, on \inserted{\handwritten{a}} $P\handwritten{<}2P$, donc, si $\underline a$ et $\underline b$ sont les idéaux divisoriels correspondant à $P$ et $2P$, on a $\underline a\handwritten{\supset}\underline b$ et $\underline a\handwritten{\ne}\underline b$ ; pour $x\handwritten{\in}\underline a-\underline b$, on a $\operatorname{div}(x)\handwritten{\geq}P$ et $\operatorname{div}(x)\handwritten{\unsure{\ngeq|\ne}}2P$, d'où $v_P(x)=1$, ce qui démontre notre assertion.

PROPOSITION 6.- \uline{Soient} $A$ \uline{un anneau de Krull}, $K$ \uline{son corps des fractions}, \uline{et}

% archive sheet number: 99
